\documentclass[12pt,a4paper]{article}

\usepackage[utf8]{inputenc}
\usepackage[T1]{fontenc}
\usepackage{lmodern}
\usepackage[margin=1in]{geometry}
\usepackage{hyperref}
\usepackage{booktabs}
\usepackage{longtable}
\usepackage{listings}
\usepackage{xcolor}
\usepackage{amssymb}
\usepackage{graphicx}
\usepackage{tikz}
\usepackage{enumitem}

\hypersetup{
    colorlinks=true,
    linkcolor=blue,
    urlcolor=blue,
}

\lstset{
    basicstyle=\ttfamily\small,
    backgroundcolor=\color{gray!10},
    frame=single,
    framerule=0pt,
    breaklines=true,
    columns=fullflexible,
    keepspaces=true,
    literate={∧}{{$\land$}}1 {∨}{{$\lor$}}1 {→}{{$\rightarrow$}}1 {▶}{{$\blacktriangleright$}}1 {✓}{{$\checkmark$}}1,
}

\title{Propositional Logic Evaluator - Report}
\author{Kyal Sin Lin Lett - st126112}
\date{28 Mar 2026}

\begin{document}

\maketitle

\newpage

%% -------------------------------------------------------------------
\section{Overview}

The task is to build an evaluator for propositional logic expressions that
supports two operators:

\begin{center}
\begin{tabular}{lll}
\toprule
Symbol & Meaning & Precedence \\
\midrule
\texttt{$\land$} & AND (conjunction) & Higher \\
\texttt{$\lor$} & OR (disjunction) & Lower \\
\bottomrule
\end{tabular}
\end{center}

The evaluator must produce \textbf{two outputs} from a single input expression:

\begin{enumerate}
    \item A \textbf{truth value} (\texttt{t} or \texttt{f})
    \item An equivalent expression in \textbf{prefix notation}
\end{enumerate}

%% -------------------------------------------------------------------
\section{Grammar}

The formal grammar of the language is:

\begin{lstlisting}
statement  ::= expr

expr       ::= expr '∨' expr       (left-associative, lower precedence)
             | expr '∧' expr       (left-associative, higher precedence)
             | 't'
             | 'f'
\end{lstlisting}

The atoms are:

\begin{itemize}
    \item \texttt{t} --- the boolean literal \emph{true}
    \item \texttt{f} --- the boolean literal \emph{false}
\end{itemize}

Precedence rule (matching the assignment specification):

\begin{quote}
$\land$ binds tighter than $\lor$, so \texttt{t $\lor$ f $\land$ f} is parsed as \texttt{t $\lor$ (f $\land$ f)}, which
evaluates to \texttt{t}.
\end{quote}

%% -------------------------------------------------------------------
\section{Type of Parser}

The parser is an \textbf{LALR(1)} parser, generated by the \texttt{sly} library's \texttt{Parser}
base class.

\textbf{Why LALR(1)?}

\begin{itemize}
    \item The grammar is unambiguous when operator precedence is declared explicitly
          (which \texttt{sly} supports via its \texttt{precedence} tuple).
    \item LALR(1) parsers are efficient (linear time) and sufficient for this simple
          grammar.
    \item It fits naturally into the existing codebase, which already uses \texttt{sly} for its
          arithmetic evaluator example.
\end{itemize}

The precedence declaration inside \texttt{PropLogicParser}:

\begin{lstlisting}[language=Python]
precedence = (
    ("left", OR),   # lower  — ∨
    ("left", AND),  # higher — ∧
)
\end{lstlisting}

Entries listed \textbf{later} have \textbf{higher} precedence. This single declaration
handles all associativity and priority conflicts automatically --- no manual
grammar rewriting needed.

%% -------------------------------------------------------------------
\section{Project Structure}

The implementation lives inside \texttt{src/prop\_evaluator/}, mirroring the provided
\texttt{src/example/} starter:

\begin{lstlisting}
src/prop_evaluator/
├── __init__.py
├── main.py                        <- PySide6 GUI entry point
└── components/
    ├── __init__.py
    ├── lexica.py                  <- Tokenizer (PropLogicLexer)
    ├── parser.py                  <- LALR(1) parser (PropLogicParser)
    ├── main.ui                    <- Qt Designer UI layout
    ├── ui.py                      <- Auto-generated from main.ui
    └── ast/
        └── prop_logic.py          <- AST node classes
\end{lstlisting}

This mirrors the starter's \texttt{src/example/} layout intentionally, keeping the same
separation of concerns:

\begin{center}
\begin{tabular}{lll}
\toprule
Layer & File & Responsibility \\
\midrule
Lexical analysis & \texttt{components/lexica.py} & Turn raw text into a token stream \\
Syntactic analysis & \texttt{components/parser.py} & Turn token stream into an AST \\
Semantic objects & \texttt{components/ast/prop\_logic.py} & Carry the logic for evaluation and translation \\
Presentation & \texttt{main.py} & GUI, wires everything together \\
\bottomrule
\end{tabular}
\end{center}

%% -------------------------------------------------------------------
\section{Step-by-Step Implementation}

\subsection{Step 1 --- AST Nodes}

\textbf{File:} \texttt{src/prop\_evaluator/components/ast/prop\_logic.py}

\textbf{Thinking:} The assignment requires two different outputs from the same parsed
expression. If we evaluate immediately during parsing (as \texttt{MyParser} in the
example does), we lose the structure needed to produce prefix notation. Instead,
we build an \textbf{Abstract Syntax Tree} during parsing and defer all computation.

This is the same reasoning that motivated \texttt{ASTParser} in the starter project.

\textbf{Design:}

\begin{itemize}
    \item \texttt{PropExpression} --- abstract base class with two abstract methods:
    \begin{itemize}
        \item \texttt{run() -> bool} --- evaluates the expression recursively
        \item \texttt{prefix() -> str} --- produces the prefix-notation string recursively
    \end{itemize}

    \item \texttt{PropTrue} / \texttt{PropFalse} --- leaf nodes, no children:

    \begin{lstlisting}[language=Python]
PropTrue().run()    # → True
PropTrue().prefix() # → "t"
    \end{lstlisting}

    \item \texttt{PropAnd(left, right)} / \texttt{PropOr(left, right)} --- internal nodes with two
          children:
    \begin{lstlisting}[language=Python]
PropAnd(PropFalse(), PropFalse()).run()    # → False
PropAnd(PropFalse(), PropFalse()).prefix() # → "∧ f f"
    \end{lstlisting}
\end{itemize}

Both \texttt{run()} and \texttt{prefix()} work by \textbf{recursively} calling the same method on
child nodes before combining results. This is the standard tree-traversal
pattern for AST evaluation.

%% -------------------------------------------------------------------

\subsection{Step 2 --- Lexer}

\textbf{File:} \texttt{src/prop\_evaluator/components/lexica.py}

\textbf{Thinking:} The lexer's job is simple: break the raw input string into a
stream of tokens that the parser can consume. For this language there are only 4
token types.

\textbf{Tokens defined:}

\begin{center}
\begin{tabular}{lll}
\toprule
Token & Pattern & Notes \\
\midrule
\texttt{AND} & $\land$ & Unicode U+2227 \\
\texttt{OR} & $\lor$ & Unicode U+2228 \\
\texttt{TRUE} & \texttt{t} & Single character \\
\texttt{FALSE} & \texttt{f} & Single character \\
\bottomrule
\end{tabular}
\end{center}

Whitespace (spaces, tabs, newlines) is silently ignored via \texttt{ignore = " \textbackslash t\textbackslash n"}.

The operators are defined \textbf{before} the atoms in the file. In \texttt{sly},
string-type token rules are matched in order of \textbf{decreasing length} (longer
patterns first), but keeping operators first is cleaner and avoids any potential
ambiguity with future tokens.

%% -------------------------------------------------------------------

\subsection{Step 3 --- Parser}

\textbf{File:} \texttt{src/prop\_evaluator/components/parser.py}

\textbf{Thinking:} The parser receives the token stream from the lexer and constructs
the AST. Each grammar rule's action creates and returns an AST node rather than
evaluating immediately.

\textbf{Grammar rules $\rightarrow$ AST construction:}

\begin{lstlisting}[language=Python]
@_("expr OR expr")
def expr(self, p) -> PropExpression:
    return PropOr(p.expr0, p.expr1)   # build a node, don't evaluate

@_("expr AND expr")
def expr(self, p) -> PropExpression:
    return PropAnd(p.expr0, p.expr1)

@_("TRUE")
def expr(self, p) -> PropExpression:
    return PropTrue()

@_("FALSE")
def expr(self, p) -> PropExpression:
    return PropFalse()
\end{lstlisting}

The precedence declaration:

\begin{lstlisting}[language=Python]
precedence = (
    ("left", OR),    # <- lower
    ("left", AND),   # <- higher
)
\end{lstlisting}

\ldots resolves all shift/reduce conflicts automatically. \texttt{sly}/\texttt{yacc}-style parsers
consult this table when they encounter an ambiguity and choose the action that
respects the declared precedence.

\textbf{Integration with the lexer:}

\begin{lstlisting}[language=Python]
tokens = PropLogicLexer.tokens
\end{lstlisting}

This single line pulls the token set from the lexer into the parser, the same
pattern used in the starter project.

%% -------------------------------------------------------------------

\subsection{Step 4 --- GUI}

\textbf{File:} \texttt{src/prop\_evaluator/main.py}

\textbf{Thinking:} The GUI layout is defined in \texttt{components/main.ui} using Qt
Designer's XML format, then compiled to \texttt{components/ui.py} via \texttt{pyside6-uic}.
\texttt{main.py} loads this compiled class (\texttt{Ui\_MainWindow}) and binds each button to
its handler --- the same pattern used in the starter project's \texttt{src/example/}.

\textbf{Layout:}

\begin{lstlisting}
[ Input: _________________________________ ]

[ t ]  [ f ]  [ ∧ ]  [ ∨ ]  [ Clear ]  [ Evaluate = ]

-----------------------------------------
Truth Value:   t
Prefix:        ∨ t ∧ f f
\end{lstlisting}

\textbf{Wiring:}

\begin{enumerate}
    \item Each atom/operator button appends its character to the input field.
    \item \texttt{Clear} resets the input and both output labels.
    \item \texttt{Evaluate} (or pressing Enter) runs the full pipeline:
    \begin{lstlisting}
input string
    → PropLogicLexer.tokenize()
    → PropLogicParser.parse()
    → AST root node
    → root.run()    → truth value label
    → root.prefix() → prefix label
    \end{lstlisting}
\end{enumerate}

%% -------------------------------------------------------------------
\section{Method of Translation and Integration}

The translation from infix expression to both outputs follows the \textbf{AST
interpreter} pattern:

\begin{lstlisting}
Input text  --[Lexer]--▶  Token stream
                              |
                          [Parser]
                              |
                           AST root
                          /        \
                   [run()]        [prefix()]
                      |                |
               bool result       prefix string
                      |                |
               truth label      prefix label
\end{lstlisting}

\textbf{Key design decision --- deferred evaluation:} The parser never computes a
value. It only \emph{builds} the tree. The two outputs are produced by two
independent traversals of the same tree after parsing completes. This cleanly
separates:

\begin{itemize}
    \item \textbf{Parsing} (syntax) from \textbf{evaluation} (semantics)
    \item The two different semantic interpretations (truth value vs.\ prefix notation)
\end{itemize}

%% -------------------------------------------------------------------
\section{Example Walkthrough}

Input: \texttt{t $\lor$ f $\land$ f}

\textbf{Lexing:}

\begin{lstlisting}
TRUE  OR  FALSE  AND  FALSE
\end{lstlisting}

\textbf{Parsing} (respecting AND > OR precedence):

\begin{lstlisting}[language=Python]
PropOr(
    PropTrue(),
    PropAnd(PropFalse(), PropFalse())
)
\end{lstlisting}

\textbf{\texttt{run()} traversal:}

\begin{lstlisting}
PropOr.run()
  → PropTrue().run()              = True
  → PropAnd(f, f).run()
      → PropFalse().run()         = False
      → PropFalse().run()         = False
      → False and False           = False
  → True or False                 = True   ✓
\end{lstlisting}

\textbf{\texttt{prefix()} traversal:}

\begin{lstlisting}
PropOr.prefix()
  → "∨ " + PropTrue().prefix() + " " + PropAnd(f,f).prefix()
  → "∨ " + "t"                  + " " + "∧ f f"
  → "∨ t ∧ f f"                                              ✓
\end{lstlisting}

\end{document}
